\documentclass[a4paper,11pt]{ctexart}

\usepackage[top=1.8cm,bottom=1.8cm,left=2.1cm,right=2.1cm]{geometry}
\usepackage{amsmath,amssymb}
\usepackage{enumitem}
\usepackage{fancyhdr}
\usepackage{lastpage}

\pagestyle{fancy}
\fancyhf{}
\renewcommand{\headrulewidth}{0pt}
\fancyfoot[C]{\small 第~\thepage~页\quad 共~\pageref{LastPage}~页}

\linespread{1.28}
\setlength{\parindent}{0pt}

\newcommand{\tr}{\operatorname{tr}}
\newcommand{\rk}{\mathrm{r}}
\newcommand{\T}{\mathrm{T}}
\newcommand{\blank}[1][2.8cm]{\underline{\hspace{#1}}}

\begin{document}

\begin{center}
  {\Large\bfseries 华中科技大学 2019～2020 学年第二学期}\\[0.4em]
  {\Large\bfseries “线性代数”考试试卷（A 卷）}\\[1.0em]
  {\normalsize
    考试方式：\underline{\ \ 闭卷\ \ }\quad
    考试日期：\underline{\ 2020.09.02AM\ }\quad
    考试时长：\underline{\ 150 分钟\ }
  }
\end{center}

\vspace{0.8em}

\textbf{一、判断题（$2\text{ 分}\times 8 = 16\text{ 分}$）正确的打“$\checkmark$”，错误的打“$\times$”}

\begin{enumerate}[label={(\quad)\ \arabic*.}, leftmargin=3.8em, itemsep=0.55em, topsep=0.6em]
  \item 设 $A$ 是 $n$ 阶对称矩阵，$B$ 为 $n$ 阶反对称矩阵，则 $AB + BA$ 是对称矩阵.

  \item 设 $A$ 为 $n$ 阶方阵，若 $A^2 = A$，则 $A = 0$ 或 $A = I$.

  \item 若方阵 $A$ 的列向量组是正交向量组，则 $A$ 为正交矩阵.

  \item 设 $A$、$B$ 都是 $n$ 阶方阵，若 $AB + A + B = 0$，则 $AB = BA$.

  \item 若 $A$ 和 $A^{\T}$ 都是 $n$ 阶方阵，则它们有相同的特征值.

  \item 设三维向量 $\beta$ 在基 $\{\alpha_1, \alpha_2, \alpha_3\}$ 下的坐标为 $(1, -1, 1)^{\T}$，则 $\beta$ 在另一组基 $\{\alpha_1, \alpha_1 + \alpha_2, \alpha_1 - \alpha_2 - \alpha_3\}$ 下的坐标为 $(4, -2, 1)^{\T}$.

  \item 二次型 $f(x, y, z) = x^2 + 2y^2 - 2z^2 - 4xy - 4yz$ 对应的矩阵的秩是 $3$.

  \item 设三阶方阵 $A$ 的特征值为 $1, -1, 2$，则行列式 $|A^2 - 3A + 5I| = 60$.
\end{enumerate}

\vspace{1.0em}

\textbf{二、填空题（$4\text{ 分}\times 5 = 20\text{ 分}$）将答案填在指定的横线上}

\begin{enumerate}[label=\arabic*., leftmargin=2em, itemsep=1.0em, topsep=0.6em]
  \item 设 $A = \begin{bmatrix} 1 & 2 & 2 \\ 2 & 1 & 2 \\ 2 & 2 & 1 \end{bmatrix}$，则 $A^5 = \blank[3.6cm]$.

  \item 方程组 $\begin{cases} x_1 + x_2 - 4x_3 + x_4 = 0 \\ x_2 - x_3 + 2x_4 = 0 \\ x_1 - x_2 - 2x_3 - 3x_4 = 0 \end{cases}$ 的通解为 $x = \blank[4.2cm]$.

  \item 设 $-2, 4$ 为矩阵 $A = \begin{bmatrix} 1 & -3 & 3 \\ 3 & a & 3 \\ 6 & -6 & b \end{bmatrix}$ 的特征值，则 $a = \blank[1.8cm]$，$b = \blank[1.8cm]$.

  \item 向量组 $\alpha_1 = \begin{bmatrix} 3 \\ 2 \\ -1 \\ 1 \end{bmatrix},\ \alpha_2 = \begin{bmatrix} 4 \\ 2 \\ 1 \\ 1 \end{bmatrix},\ \alpha_3 = \begin{bmatrix} -2 \\ -2 \\ 3 \\ -1 \end{bmatrix}$ 的秩为 $\blank[2.2cm]$.

  \item 设矩阵 $A = (\alpha_1, \alpha_2, \alpha_3, \alpha_4)$，其中 $\alpha_1, \alpha_2, \alpha_3, \alpha_4$ 都是四维列向量，$P$ 为四阶矩阵，若 $AP = (\alpha_1, \alpha_2, \alpha_3, 2\alpha_1 + \alpha_4)$，则 $P = \blank[3.6cm]$.
\end{enumerate}

\vspace{1.0em}

\textbf{三、计算 $n$ 阶行列式的值（8 分）}
\[
  D_n = \begin{vmatrix}
    a+x & b & b & \cdots & b & b \\
    a & a+x & 0 & \cdots & 0 & 0 \\
    a & 0 & a+x & \cdots & 0 & 0 \\
    \vdots & \vdots & \vdots & \ddots & \vdots & \vdots \\
    a & 0 & 0 & \cdots & a+x & 0 \\
    a & 0 & 0 & \cdots & 0 & a+x
  \end{vmatrix} \text{，其中 } a + x \neq 0.
\]

\vspace{0.8em}

\textbf{四、（12 分）}已知向量 $\alpha_1 = \begin{bmatrix} 1 \\ a \\ 1 \end{bmatrix},\ \alpha_2 = \begin{bmatrix} 1 \\ 1 \\ a \end{bmatrix},\ \alpha_3 = \begin{bmatrix} a \\ 1 \\ 1 \end{bmatrix},\ \beta = \begin{bmatrix} -2 \\ a - 3 \\ -2 \end{bmatrix}$，问 $a$ 为何值时：
\begin{enumerate}[label=(\arabic*), leftmargin=2.4em, itemsep=0.25em, topsep=0.3em]
  \item $\beta$ 可由 $\{\alpha_1, \alpha_2, \alpha_3\}$ 唯一线性表示；
  \item $\beta$ 不能由 $\{\alpha_1, \alpha_2, \alpha_3\}$ 线性表示；
  \item $\beta$ 可由 $\{\alpha_1, \alpha_2, \alpha_3\}$ 线性表示但不唯一，并写出此线性表示式.
\end{enumerate}

\vspace{0.8em}

\textbf{五、（12 分）}用正交变换化二次型 $f(x_1, x_2, x_3) = 2x_1 x_2 + 2x_1 x_3 - 2x_2 x_3$ 为标准形.

\vspace{0.8em}

\textbf{六、（12 分）}设 $A = (\alpha_1, \alpha_2, \alpha_3, \alpha_4)$ 为四阶方阵，$\alpha_1, \alpha_2, \alpha_3, \alpha_4$ 都是 $4$ 维列向量，其中 $\alpha_2, \alpha_3, \alpha_4$ 线性无关，且 $\alpha_1 = \alpha_2 - 2\alpha_3$，$\beta = \alpha_1 + \alpha_2 - \alpha_3 - \alpha_4$，求方程组 $AX = \beta$ 的通解.

\vspace{0.8em}

\textbf{七、（10 分）}设 $A$ 为实对称矩阵，且 $(A - I)(A^2 + 2A + 5I) = 0$. 证明 $A = I$.

\vspace{0.8em}

\textbf{八、（10 分）}设 $A = \begin{bmatrix}
  a_{11} & a_{12} & \cdots & a_{1n} \\
  0 & a_{22} & \cdots & a_{2n} \\
  \vdots & \vdots & \ddots & \vdots \\
  0 & 0 & \cdots & a_{nn}
\end{bmatrix}$ 为上三角矩阵，证明
\begin{enumerate}[label=\arabic*., leftmargin=2.4em, itemsep=0.25em, topsep=0.3em]
  \item 当 $a_{ii} \neq a_{jj}\ (i \neq j)$ 时，$A$ 可对角化.
  \item 当 $a_{11} = a_{22} = \cdots = a_{nn}$ 时，若存在 $a_{ij} \neq 0\ (i < j)$，则 $A$ 不能相似于对角阵.
\end{enumerate}

\end{document}
